\documentclass[11pt]{article}

\usepackage[margin=1in]{geometry}
\usepackage{amsmath,amssymb,amsthm,mathtools}
\usepackage{enumitem}
\usepackage[colorlinks=true,linkcolor=blue,citecolor=blue,urlcolor=blue]{hyperref}

\newtheorem{conjecture}{Conjecture}
\newtheorem{problem}{Problem}
\newtheorem{lemma}{Lemma}
\newtheorem{proposition}{Proposition}
\newtheorem{corollary}{Corollary}
\newtheorem{heuristic}{Heuristic}
\newtheorem{definition}{Definition}
\newtheorem{example}{Example}

\DeclareMathOperator{\supp}{supp}
\DeclareMathOperator{\TC}{TC}
\DeclareMathOperator{\Ent}{H}
\DeclareMathOperator{\E}{\mathbb{E}}
\DeclareMathOperator{\Prb}{\mathbb{P}}

\title{An Entropy--Transport Attack on Frankl's Union-Closed Sets Conjecture}
\author{Open Conjecture Formalizations}
\date{\today}

\begin{document}

\maketitle

\begin{abstract}
This note develops a possible route toward Frankl's union-closed sets
conjecture by trying to combine two proof paradigms that have historically
stalled separately: entropy growth under union, and pairing/matching
arguments.  The organizing idea is that the entropy method fails precisely
when the union map has large collision/dependence cost, while pairing methods
fail precisely when attempted injections collide.  We propose treating these
as the same obstruction.  Either a well-chosen coupling of two random members
of the family produces an entropy contradiction, or the dual obstruction to
such a coupling should produce a Hall-type certificate rigid enough to force
a Frankl element.
\end{abstract}

\section{The Basic Setup}

Let $\mathcal{F}$ be a finite union-closed family of finite sets.  After
replacing the ambient set by $\bigcup_{A \in \mathcal{F}} A$, write the ground
set as $[n]$.  Equivalently,
\[
  \mathcal{F} \subseteq \{0,1\}^n
\]
is closed under coordinatewise OR.

Let $X$ be uniformly distributed on $\mathcal{F}$, and write
\[
  p_i = \Prb[X_i = 1].
\]
Frankl's conjecture is:
\[
  \max_i p_i \geq \frac12,
\]
provided $\mathcal{F}$ has at least one nonempty member.

\begin{conjecture}[Frankl]
If $\mathcal{F}$ is finite, nontrivial, and union-closed, then some element of
the ground set belongs to at least half of the members of $\mathcal{F}$.
\end{conjecture}

For this note, assume toward contradiction that $\mathcal{F}$ is a
counterexample:
\[
  0 < p_i < \frac12 \qquad \text{for every active coordinate } i.
\]

\section{The Join-Semilattice View}

The family $\mathcal{F}$ is a finite join-semilattice with join operation
\[
  A \vee B = A \cup B.
\]
Each coordinate $i$ gives a join-homomorphism
\[
  \chi_i : \mathcal{F} \to \mathbf{2},
  \qquad
  \chi_i(A) = 1 \Longleftrightarrow i \in A,
\]
because
\[
  \chi_i(A \vee B) = \chi_i(A) \vee \chi_i(B).
\]
The maps $\chi_i$ jointly separate points if the ground set has been reduced
by identifying indistinguishable elements.

Thus Frankl can be read as a sharp fiber-size statement:
\[
  \text{one separating join-homomorphism }
  \chi_i : \mathcal{F} \to \mathbf{2}
  \text{ has }
  |\chi_i^{-1}(1)| \geq |\mathcal{F}|/2.
\]
The categorical viewpoint is useful because it isolates the structure:
\[
  \text{finite join-semilattice}
  \quad + \quad
  \text{separating coordinates to } \mathbf{2}
  \quad + \quad
  \text{uniform counting measure}.
\]
The hard part is the last phrase.  Pure structure-preserving maps do not, by
themselves, remember the sharp counting threshold.

\section{Known Quantitative Barriers}

\subsection{The iid Entropy Barrier}

Let $X,Y$ be independent uniform samples from $\mathcal{F}$, and set
\[
  Z = X \vee Y.
\]
Since $\mathcal{F}$ is union-closed, $Z$ is supported on $\mathcal{F}$, so
\[
  \Ent(Z) \leq \log |\mathcal{F}| = \Ent(X).
\]
On coordinate $i$,
\[
  q_i := \Prb[Z_i = 1] = 2p_i - p_i^2.
\]
The coordinate entropy gain is
\[
  h(q_i) - h(p_i),
\]
where $h(t)=-t\log t-(1-t)\log(1-t)$.

This gain is positive as long as
\[
  p_i < 2p_i - p_i^2 < 1-p_i,
\]
which gives
\[
  p_i < \frac{3-\sqrt{5}}{2}.
\]
This is the natural barrier faced by iid OR-sampling: above that threshold,
OR pushes the coordinate past the entropy peak at $1/2$.

\subsection{The Gilmer--AHS Entropy Hinge}

Gilmer's breakthrough isolates a one-coordinate conditional entropy
inequality.  After revealing a history $C$, write
\[
  p_C=\Prb[X_i=1\mid C],
  \qquad
  x_C=1-p_C.
\]
For independent histories $C,C'$, the probability that the OR bit is zero is
$x_Cx_{C'}$, so the conditional entropy contribution is
\[
  \E_{C,C'} h(x_Cx_{C'}).
\]
Alweiss--Huang--Sellke reduce the sharp version of Gilmer's lemma to the
following optimization over probability measures $\mu$ on $[0,1]$ with mean
$\phi$:
\[
  F(\mu)=\E_{x,y\sim\mu\times\mu} h(xy)-\E_{x\sim\mu}h(x).
\]
They prove that the minimum is attained on a measure supported on
$\{0,x\}$, so the hinge is the one-variable inequality
\[
  \phi h(x^2)\ge xh(x)\qquad (x\in[\phi,1]).
\]
For
\[
  \phi=\frac{\sqrt5-1}{2},
\]
this inequality holds, with equality at $x=\phi$ and $x=1$.  Therefore the
iid entropy method reaches the constant
\[
  1-\phi=\frac{3-\sqrt5}{2}.
\]

The companion script
\[
  \texttt{gilmer\_ahs\_verification.py}
\]
reproduces the finite numerical part of this hinge proof.  In particular,
\[
  \texttt{--verify-ahs}
\]
checks the $I_1$ convexity table, the $I_2$ monotone chain, and the final
$I_3$ analytic margin from the Alweiss--Huang--Sellke appendix.  This is an
audit scaffold rather than a Lean certificate, but the checks are deliberately
separated into the finite inequalities that should eventually be replaced by
interval arithmetic or Lean proofs.
The same script can also write
\[
  \texttt{ahs\_hinge\_certificate.json}
\]
with the rational grid points and high-precision margins for the finite
$I_1$ and $I_2$ checks.  The payload now also includes mpmath interval
enclosures and conservative rational floors for each finite row.  This is the
bridge format for the Lean checker: it separates the finite row checking from
the still-open task of formalizing the transcendental interval enclosures.

This also marks the exact point where our technique must diverge from iid
entropy.  The product measure $\mu\times\mu$ is forced by independent samples.
A first, weak coupled analogue would replace it by a transport plan $\kappa$
with the same marginals, and ask whether one can make
\[
  \E_{(x,y)\sim\kappa}h(xy)-\E_{x\sim\mu}h(x)
\]
positive even in regimes where the product plan has already crossed the
golden-ratio barrier.  But the genuine coupled problem is stronger still:
after choosing a pair of histories with conditional zero probabilities
$x$ and $y$, the conditional coupling of the two fibers may choose a
zero-zero mass
\[
  t\in[\max(0,x+y-1),\min(x,y)]
\]
rather than being forced to use $t=xy$.  Thus the coupled hinge should be an
optimization of $\E h(t)$ over both a history transport plan and Frechet
feasible conditional kernels.

At the raw one-coordinate level this is plausible:
the script's
\[
  \texttt{--coupled-coordinate-probe}
\]
shows that centering an OR coordinate at entropy $1/2$ keeps positive
coordinate entropy gain for every $p<1/2$, while iid OR becomes negative
above $(3-\sqrt5)/2$.  The remaining problem is global: such coordinatewise
transport must be made compatible across all histories and all coordinates,
and its total-correlation cost must be paid by the fiber/Hall certificates
developed below.

\begin{proposition}[One-coordinate coupled OR window]
Let $X_i$ and $Y_i$ be Bernoulli random variables with the same marginal
$p=\Prb[X_i=1]=\Prb[Y_i=1]$, coupled arbitrarily.  If
\[
  r=\Prb[X_i=1,\ Y_i=1],
\]
then the OR marginal is
\[
  q=\Prb[X_i\vee Y_i=1]=2p-r.
\]
As $r$ ranges over the Frechet interval
\[
  \max(0,2p-1)\le r\le p,
\]
the attainable OR marginals are exactly
\[
  p\le q\le \min(2p,1).
\]
Consequently every coordinate with $1/4\le p\le 1/2$ can be centered at
$q=1/2$, by taking $r=2p-1/2$; and every coordinate with $0<p<1/4$ can still
be moved upward inside the entropy-increasing region, by taking $r=0$ and
$q=2p$.
\end{proposition}

\begin{proof}
The identity $q=2p-r$ is inclusion--exclusion.  The bounds on $r$ are the
two-marginal Frechet bounds.  Substituting $r=2p-q$ gives
$p\le q\le\min(2p,1)$.  The two displayed choices of $r$ are then immediate.
The algebraic feasibility part of this proposition is formalized in the Lean
file \texttt{CoupledOr.lean}.
\end{proof}

The companion script also has a
\[
  \texttt{--one-coordinate-frontier}
\]
mode which prints this interval next to the iid value.  This makes the
strategy's intended constant improvement precise: at the one-coordinate level
there is no golden-ratio barrier left.  The whole difficulty has moved into
the existence of one global coupling whose coordinate shadows satisfy these
local demands simultaneously.

\begin{problem}[Frechet-kernel replacement for the AHS hinge]
Find structural hypotheses, coming from union-closure rather than from a
single coordinate alone, under which the following coupled hinge is positive.
Given a distribution $\mu$ of conditional zero probabilities, choose a
transport plan $\kappa$ with both marginals $\mu$, and for each pair $(x,y)$
choose
\[
  t(x,y)\in[\max(0,x+y-1),\min(x,y)].
\]
When can one force
\[
  \E_{(x,y)\sim\kappa} h(t(x,y))>\E_{x\sim\mu}h(x),
\]
after subtracting the total-correlation cost needed to realize these
coordinatewise kernels simultaneously?
\end{problem}

This is the precise formal target suggested by the computation.  The AHS
paper proves that the product specialization $t=xy$ cannot pass the
golden-ratio constant.  Our coupled program asks whether union-closure supplies
enough transport structure to choose better $t$ values without paying more
dependence than the marginal entropy gain earns.

\subsection{Average-Size Bounds Lose Dimension}

Average-size methods estimate
\[
  \sum_i p_i = \E[|X|].
\]
But Frankl requires one $p_i \geq 1/2$.  Any lower bound on the sum loses force
when the effective number of coordinates is large.

\subsection{Pairing Arguments Collide}

If $\{i\} \in \mathcal{F}$, then the map
\[
  A \mapsto A \cup \{i\}
\]
injects members omitting $i$ into members containing $i$, proving that $i$
appears in at least half the family.

For a larger set $S \in \mathcal{F}$, the natural map
\[
  A \mapsto A \cup S
\]
can collapse many different $A$ to the same image.  This is the basic local
failure of pairing proofs.

\section{Coupled OR Transport}

The iid choice of $(X,Y)$ is only one coupling of the uniform distribution on
$\mathcal{F}$ with itself.  Let $\Pi(\mathcal{F})$ be the set of probability
measures $\pi$ on $\mathcal{F} \times \mathcal{F}$ whose two marginals are both
uniform.

For $\pi \in \Pi(\mathcal{F})$, let $(X,Y) \sim \pi$ and
\[
  Z = X \vee Y.
\]
Define
\[
  q_i(\pi) = \Prb_\pi[Z_i=1].
\]
Since $Z$ is supported on $\mathcal{F}$,
\[
  \Ent(Z) \leq \Ent(X)=\log|\mathcal{F}|.
\]
Now decompose entropy into coordinate entropy minus total correlation:
\[
  \Ent(W)
  =
  \sum_i \Ent(W_i) - \TC(W),
\]
where
\[
  \TC(W) = \sum_i \Ent(W_i) - \Ent(W).
\]
Thus
\[
  \Ent(Z)-\Ent(X)
  =
  \sum_i \bigl(h(q_i(\pi))-h(p_i)\bigr)
  -
  \bigl(\TC(Z)-\TC(X)\bigr).
\]

\begin{definition}[Entropy advantage]
For a coupling $\pi \in \Pi(\mathcal{F})$, define
\[
  \Phi(\pi)
  =
  \sum_i \bigl(h(q_i(\pi))-h(p_i)\bigr)
  -
  \bigl(\TC(Z)-\TC(X)\bigr).
\]
\end{definition}

Since $Z$ is supported on $\mathcal{F}$, every coupling satisfies
\[
  \Phi(\pi) = \Ent(Z)-\Ent(X) \leq 0.
\]
Therefore:

\begin{proposition}[Entropy contradiction criterion]
If, under the assumption $p_i<1/2$ for all $i$, there exists
$\pi\in\Pi(\mathcal{F})$ with $\Phi(\pi)>0$, then $\mathcal{F}$ is not a
counterexample.
\end{proposition}

The difficulty is to choose $\pi$.  The iid coupling often makes
$q_i(\pi)=2p_i-p_i^2$ too large for coordinates near $1/2$.  The diagonal
coupling makes $q_i(\pi)=p_i$ and gives no coordinate gain.  We want a
positively correlated but not fully diagonal coupling:
\[
  p_i < q_i(\pi) < 1-p_i
  \qquad \text{for all } i.
\]

Equivalently, since
\[
  q_i(\pi)=2p_i-r_i(\pi),
  \qquad
  r_i(\pi)=\Prb_\pi[X_i=Y_i=1],
\]
we want
\[
  3p_i-1 < r_i(\pi) < p_i.
\]
For coordinates close to $1/2$, this asks for very strong positive
correlation, but still leaves room below the diagonal.

\section{The Linear Feasibility Problem and Its Dual}

Fix a set of coordinates $I$ and desired targets $\tau_i$ for $i\in I$.
For example, $\tau_i=1/2$ asks to center coordinate $i$ at the entropy peak.
Let
\[
  c_i(A,B)=\mathbf{1}_{i\in A\cup B}.
\]
The centering problem is the following transportation feasibility problem:
\[
\begin{array}{rll}
  \text{find} & \pi_{A,B}\ge 0 & (A,B\in\mathcal{F})\\[2mm]
  \text{such that}
    & \displaystyle \sum_B \pi_{A,B}=1/m & (A\in\mathcal{F}),\\[3mm]
    & \displaystyle \sum_A \pi_{A,B}=1/m & (B\in\mathcal{F}),\\[3mm]
    & \displaystyle \sum_{A,B} c_i(A,B)\pi_{A,B}=\tau_i
      & (i\in I),
\end{array}
\]
where $m=|\mathcal{F}|$.

By Farkas' lemma, this system is infeasible if and only if there exist row
potentials $u_A$, column potentials $v_B$, and coordinate weights $\lambda_i$
such that
\[
  u_A+v_B+\sum_{i\in I}\lambda_i c_i(A,B)\ge 0
  \qquad\text{for all }A,B\in\mathcal{F},
\]
but
\[
  \frac1m\sum_A u_A+\frac1m\sum_B v_B+\sum_{i\in I}\lambda_i\tau_i<0.
\]
Thus any failure to center a collection of coordinates has a very concrete
dual certificate: a weighted inequality on the OR-incidence matrix that is
nonnegative on every pair $(A,B)$ but negative on the demanded marginals.

For the window problem one replaces the equality constraints by
\[
  p_i+\gamma \leq \sum_{A,B}c_i(A,B)\pi_{A,B}
  \leq 1-p_i-\gamma
\]
and maximizes $\gamma$.  The dual then has two nonnegative coordinate weights
for each $i$, one for the lower wall and one for the upper wall.  Boundary
examples with $p_i=1/2$ are already explained by this dual: the diagonal
coupling gives $q_i=p_i$, while the upper wall is also $1-p_i=p_i$, so no
positive margin is possible.

\section{One Coordinate Can Be Understood Exactly}

The first sanity check is the one-coordinate problem.  It does not use
union-closure at all.

\begin{proposition}[Single-coordinate OR range]
Fix a coordinate $i$, and let
\[
  R_i=\{A\in\mathcal{F}:i\in A\},\qquad
  L_i=\{A\in\mathcal{F}:i\notin A\}.
\]
Write $r=|R_i|$, $\ell=|L_i|$, $m=r+\ell$, and $p_i=r/m$.  As $\pi$ ranges
over uniform couplings of two uniform samples from $\mathcal{F}$, the possible
values of
\[
  q_i(\pi)=\Prb_\pi[i\in X\cup Y]
\]
fill the interval
\[
  p_i \le q_i(\pi)\le \min(1,2p_i).
\]
\end{proposition}

\begin{proof}
By the Birkhoff-von Neumann theorem, it is enough to consider permutation
couplings and then take convex combinations.  For a permutation coupling
$Y=\sigma(X)$, every row $A\in R_i$ contributes to $q_i$.  In addition, a row
$A\in L_i$ contributes precisely when $\sigma(A)\in R_i$.  If $t$ denotes the
number of rows of $L_i$ mapped into columns of $R_i$, then
\[
  q_i=\frac{r+t}{m}.
\]
The possible values of $t$ range from $0$ to $\min(\ell,r)$: assign as many
columns of $R_i$ as desired to rows of $L_i$, and fill the remaining rows and
columns bijectively.  Hence the permutation values range from $r/m$ to
$(r+\min(\ell,r))/m=\min(1,2r/m)$, and convex combinations fill the interval.
\end{proof}

Consequently, for a strictly light coordinate $0<p_i<1/2$ the window
\[
  p_i<q_i<1-p_i
\]
is always possible for that coordinate alone.  However, exact centering
$q_i=1/2$ is possible only when
\[
  p_i\ge \frac14.
\]
This corrects the naive ``center every light coordinate'' guess: very rare
coordinates can be moved upward, but not all the way to the entropy peak in a
single OR step.

\section{The Assignment-Dual Form of Critical Centering}

The computational evidence below suggests that critical light coordinates can
often be centered simultaneously.  It is useful to put that claim into a form
that has no hidden analysis.

Fix a coordinate set $I$.  For a permutation $\sigma$ of $\mathcal{F}$, write
\[
  q(\sigma)_i
  =
  \frac1m
  \left|\{A\in\mathcal{F}:i\in A\cup\sigma(A)\}\right|
  \qquad (i\in I).
\]
By Birkhoff-von Neumann, the possible vectors $q_I(\pi)$ over all uniform
self-couplings are exactly the convex hull of the permutation vectors
$q(\sigma)$.

\begin{proposition}[Assignment dual for critical centering]
The center vector $(1/2,\dots,1/2)\in\mathbb{R}^{I}$ belongs to this convex
hull if and only if, for every weight vector $\lambda\in\mathbb{R}^{I}$,
\[
  \max_{\sigma\in S_{\mathcal{F}}}
  \frac1m
  \sum_{A\in\mathcal{F}}
  \sum_{i\in I}
  \lambda_i\mathbf{1}_{i\in A\cup\sigma(A)}
  \geq
  \frac12\sum_{i\in I}\lambda_i.
\]
\end{proposition}

\begin{proof}
This is the separating-hyperplane theorem applied to the finite polytope
$\operatorname{conv}\{q(\sigma):\sigma\in S_{\mathcal{F}}\}$.  If the center
is outside the polytope, some linear functional $\lambda$ is larger at the
center than on every permutation vector.  Conversely, such a $\lambda$ is
exactly a separating functional.
\end{proof}

Thus critical centering is equivalent to a weighted assignment problem.  For
fixed $\lambda$, define
\[
  w_\lambda(A,B)
  =
  \sum_{i\in I}\lambda_i\mathbf{1}_{i\in A\cup B}.
\]
The desired inequality is
\[
  \max_{\sigma}
  \sum_A w_\lambda(A,\sigma(A))
  \geq
  \frac{m}{2}\sum_i\lambda_i.
\]
The weight matrix has special join form:
\[
  w_\lambda(A,B)
  =
  \sum_{i\in A\cap I}\lambda_i
  +
  \sum_{i\in I\setminus A}\lambda_i\mathbf{1}_{i\in B}.
\]
So the row contribution already present in $A$ is fixed; only the missing
positive or negative coordinate weights are assigned to columns.

This exposes a thinner quantitative barrier.  The diagonal permutation gives
$q_i=p_i$, and therefore proves the dual inequality whenever all coordinate
weights are nonpositive, since $\lambda_i p_i>\lambda_i/2$ when
$\lambda_i<0$ and $p_i<1/2$.  The average over all permutations gives
\[
  \mathbb{E}_\sigma q(\sigma)_i=2p_i-p_i^2,
\]
which is at least $1/2$ when
\[
  p_i\geq 1-\frac1{\sqrt{2}}\approx 0.2929.
\]
Therefore nonnegative weights on the upper part of the critical range are
handled by the random-permutation average.  The remaining genuinely hard zone
is narrow but important:
\[
  \frac14\leq p_i<1-\frac1{\sqrt{2}},
\]
especially when these rare positive coordinates must be centered while
negative coordinates are kept from being pushed upward.  Thus the real
assignment problem is a signed compromise between diagonal behavior and
outward-pushing behavior.  A proof of the assignment dual would need to show
that the maximizing permutation beats the random average in precisely this
rare-coordinate regime.

\begin{problem}[Weighted critical assignment inequality]
Let $\mathcal{F}$ be union-closed, let $I=\{i:1/4\le p_i<1/2\}$, and let
$\lambda\in\mathbb{R}^{I}$.  Prove
\[
  \max_{\sigma\in S_{\mathcal{F}}}
  \frac1m
  \sum_A w_\lambda(A,\sigma(A))
  \geq
  \frac12\sum_i\lambda_i,
\]
or classify exactly which signed weights can fail.  Any genuine failure gives
a concrete separating certificate for the critical-centering conjecture.
\end{problem}

\subsection{The Assignment LP Dual}

For a fixed $\lambda$, write
\[
  s_\lambda(A)=\sum_{i\in A\cap I}\lambda_i,
  \qquad
  g_\lambda(A,B)=\sum_{i\in I\setminus A}\lambda_i\mathbf{1}_{i\in B}.
\]
Then
\[
  w_\lambda(A,B)=s_\lambda(A)+g_\lambda(A,B).
\]
The row term is independent of the assignment:
\[
  \sum_{A\in\mathcal{F}}s_\lambda(A)
  =
  m\sum_{i\in I}\lambda_i p_i.
\]
Therefore the weighted critical assignment inequality is equivalent to
\[
  \max_{\sigma}
  \sum_A g_\lambda(A,\sigma(A))
  \geq
  m\sum_{i\in I}\lambda_i\left(\frac12-p_i\right).
\]
Set
\[
  d_i=\frac12-p_i>0.
\]
The normalized problem is:
\[
  \max_{\sigma}
  \sum_A
  \sum_{i\in I\setminus A}\lambda_i\mathbf{1}_{i\in\sigma(A)}
  \geq
  m\sum_i\lambda_i d_i.
\]

This form is conceptually useful.  The positive part of $\lambda$ asks rows
missing $i$ to be matched into columns containing $i$.  The negative part asks
rows missing $i$ to avoid such columns.  Since the right side is also signed,
large negative weights are usually harmless; they make the demanded lower
bound easier, and the diagonal assignment realizes zero gain on every missing
coordinate.

Now pass to the bipartite assignment LP:
\[
\begin{array}{rll}
  \text{maximize}
    & \displaystyle\sum_{A,B}g_\lambda(A,B)x_{A,B}\\[2mm]
  \text{subject to}
    & \displaystyle\sum_B x_{A,B}=1 & (A\in\mathcal{F}),\\[2mm]
    & \displaystyle\sum_A x_{A,B}=1 & (B\in\mathcal{F}),\\[2mm]
    & x_{A,B}\geq 0.
\end{array}
\]
The Birkhoff-von Neumann theorem says its optimum is attained by a
permutation.  Its dual is
\[
\begin{array}{rll}
  \text{minimize}
    & \displaystyle\sum_A \alpha_A+\sum_B\beta_B\\[2mm]
  \text{subject to}
    & \alpha_A+\beta_B\geq g_\lambda(A,B)
      & (A,B\in\mathcal{F}).
\end{array}
\]
Thus failure of critical centering is equivalent to the existence of
$\lambda$, row potentials $\alpha_A$, and column potentials $\beta_B$ such
that
\[
  \alpha_A+\beta_B
  \geq
  \sum_{i\in I\setminus A}\lambda_i\mathbf{1}_{i\in B}
  \qquad(A,B\in\mathcal{F}),
\]
but
\[
  \sum_A\alpha_A+\sum_B\beta_B
  <
  m\sum_i\lambda_i d_i.
\]
This is the finite certificate we should try to rule out.

\begin{problem}[Potential obstruction]
Classify triples $(\lambda,\alpha,\beta)$ satisfying the dual inequalities
with objective below $m\sum_i\lambda_i d_i$.  Show that, after quotienting
clones and deleting dominated critical coordinates, no such triple exists.
\end{problem}

\subsection{The Remaining Obstruction Cone}

Two elementary assignments already rule out most possible signs.

First, the diagonal assignment $\sigma(A)=A$ gives
\[
  \sum_A g_\lambda(A,A)=0.
\]
Hence any failing $\lambda$ must satisfy
\[
  \sum_i\lambda_i d_i>0.
\]
Second, averaging over a uniformly random permutation gives
\[
  \mathbb{E}_\sigma
  \sum_A g_\lambda(A,\sigma(A))
  =
  m\sum_i\lambda_i p_i(1-p_i).
\]
Therefore any failing $\lambda$ must also satisfy
\[
  \sum_i\lambda_i p_i(1-p_i)
  <
  \sum_i\lambda_i d_i.
\]
Equivalently, with
\[
  e_i=p_i(1-p_i)-d_i
      =2p_i-p_i^2-\frac12,
\]
a failing $\lambda$ must lie in the cone
\[
  \sum_i\lambda_i d_i>0,
  \qquad
  \sum_i\lambda_i e_i<0.
\]
For $p_i\geq 1-1/\sqrt{2}$, one has $e_i\geq0$.  Thus positive mass on
high-critical coordinates helps the random-permutation bound rather than
hurting it.  The only positive weights that can drive the second inequality
must concentrate on the narrow low-critical range
\[
  \frac14\leq p_i<1-\frac1{\sqrt{2}}.
\]

\begin{lemma}[Necessary cone for a dual obstruction]
If $\lambda$ violates the weighted critical assignment inequality, then
\[
  \sum_i\lambda_i\left(\frac12-p_i\right)>0
  \quad\text{and}\quad
  \sum_i\lambda_i\left(2p_i-p_i^2-\frac12\right)<0.
\]
\end{lemma}

\begin{proof}
The first inequality follows because the diagonal assignment gives left-hand
side $0$ in the normalized problem.  The second follows because the optimum is
at least the average value of the normalized objective over all permutations.
\end{proof}

This cone is a useful target for both proof and search.  In particular, a
face-stripped obstruction cannot be all-negative, cannot be supported only on
high-critical positive coordinates, and cannot be explained by coordinate
order.  It must be a genuinely signed tradeoff between rare positive
coordinates and compensating negative coordinates.

\subsection{Single-Positive Directions}

The one-coordinate part of the remaining cone is completely harmless.

\begin{proposition}[One-coordinate dual inequality]
Let $I=\{i\}$ and $1/4\leq p_i<1/2$.  Then the weighted critical assignment
inequality holds for every $\lambda_i\in\mathbb{R}$.
\end{proposition}

\begin{proof}
If $\lambda_i\geq0$, choose a permutation coupling maximizing $q_i$.  By the
single-coordinate range, this gives $q_i=2p_i$, so
\[
  \lambda_i q_i-\lambda_i/2
  =
  \lambda_i(2p_i-1/2)\geq0.
\]
If $\lambda_i<0$, choose the diagonal coupling, giving $q_i=p_i$.  Then
\[
  \lambda_i q_i-\lambda_i/2
  =
  \lambda_i(p_i-1/2)\geq0,
\]
because both factors are negative.
\end{proof}

The first genuinely mixed case is therefore a positive low-critical
coordinate compensated by one or more negative coordinates.  There is a clean
matching condition that resolves many such directions.

Let $i$ be the positive coordinate and let $J$ be the set of negative
coordinates.  Partition $\mathcal{F}$ by the trace on $J$: two sets lie in the
same block if they contain exactly the same coordinates from $J$.

\begin{lemma}[Protected one-positive boosting]
Assume $1/4\leq p_i<1/2$.  Suppose that for every $J$-trace block
$\mathcal{B}\subseteq\mathcal{F}$,
\[
  |\{A\in\mathcal{B}:i\in A\}|
  \leq
  \frac{|\mathcal{B}|}{2}.
\]
Then the weighted critical assignment inequality holds for every signed
weight vector with $\lambda_i>0$ and $\lambda_j\leq0$ for all $j\in J$.
\end{lemma}

\begin{proof}
Inside each $J$-trace block $\mathcal{B}$, choose a bijection
$\sigma_{\mathcal{B}}:\mathcal{B}\to\mathcal{B}$ such that every column
containing $i$ is assigned to a row omitting $i$.  This is possible because
the number of $i$-columns in the block is at most the number of rows omitting
$i$ in the same block.  Combining these bijections gives a permutation
$\sigma$ of $\mathcal{F}$.

Since $\sigma$ preserves every $J$-trace block, it does not increase any
negative coordinate: $q_j(\sigma)=p_j$ for each $j\in J$.  Since every
$i$-column is assigned to an $i$-missing row, the extra contribution to
$q_i$ is $p_i$, so $q_i(\sigma)=2p_i$.  Hence
\[
  \sum_k\lambda_k q_k(\sigma)-\frac12\sum_k\lambda_k
  =
  \lambda_i(2p_i-1/2)
  +
  \sum_{j\in J}\lambda_j(p_j-1/2)
  \geq0.
\]
The first term is nonnegative because $p_i\geq1/4$, and each summand in the
second term is nonnegative because $\lambda_j\leq0$ and $p_j<1/2$.
\end{proof}

Thus any obstruction not handled by this lemma must have the positive
coordinate too dense inside some trace block of the negative coordinates.
This is a local obstruction, not a global entropy phenomenon.

There is a more flexible version that allows dense blocks as long as the
negative weights leave enough slack.

\begin{lemma}[Trace-preserving partial boost]
Let $i$ be the unique coordinate with positive weight, and let
$J$ be the set of coordinates with negative weights.  Partition
$\mathcal{F}$ into blocks by the $J$-trace.  For a block $\mathcal{B}$, write
\[
  r_{\mathcal{B}}=|\{A\in\mathcal{B}:i\in A\}|,
  \qquad
  \ell_{\mathcal{B}}=|\{A\in\mathcal{B}:i\notin A\}|.
\]
Set
\[
  \Delta_{i|J}
  =
  \frac1m\sum_{\mathcal{B}}\min(r_{\mathcal{B}},\ell_{\mathcal{B}}).
\]
Then there is a permutation coupling such that
\[
  q_j=p_j\quad(j\in J),
  \qquad
  q_i=p_i+\Delta_{i|J}.
\]
Consequently the weighted critical assignment inequality holds whenever
\[
  \Delta_{i|J}
  \geq
  \left(\frac12-p_i\right)
  +
  \frac1{\lambda_i}\sum_{j\in J}\lambda_j\left(\frac12-p_j\right).
\]
\end{lemma}

\begin{proof}
Inside each $J$-trace block $\mathcal{B}$, match as many $i$-containing
columns as possible to $i$-omitting rows.  This creates exactly
$\min(r_{\mathcal{B}},\ell_{\mathcal{B}})$ new OR occurrences of $i$ inside
the block.  Complete the partial matching arbitrarily to a bijection of
$\mathcal{B}$.  Combining these bijections over all blocks gives a
permutation of $\mathcal{F}$ preserving every $J$-trace, hence preserving
$q_j=p_j$ for all $j\in J$, and increasing $q_i$ by $\Delta_{i|J}$.

For this coupling,
\[
  \sum_k\lambda_k(q_k-1/2)
  =
  \lambda_i\left(p_i+\Delta_{i|J}-\frac12\right)
  +
  \sum_{j\in J}\lambda_j\left(p_j-\frac12\right).
\]
The displayed lower bound on $\Delta_{i|J}$ is exactly the condition that
this quantity is nonnegative.
\end{proof}

The protected-block lemma is the special case where
$r_{\mathcal{B}}\leq \ell_{\mathcal{B}}$ in every block, so
$\Delta_{i|J}=p_i$.

\subsection{Multi-Positive Trace Blocks}

Now let $P=\{i:\lambda_i>0\}$ and $J=\{j:\lambda_j<0\}$.  The natural
generalization is to preserve the $J$-trace while optimizing the positive
coordinates inside each trace block.

For a $J$-trace block $\mathcal{B}$, define the positive gain weight
\[
  g^+_\lambda(A,B)
  =
  \sum_{i\in P\setminus A}\lambda_i\mathbf{1}_{i\in B}.
\]
Let
\[
  M_{\mathcal{B}}(\lambda)
  =
  \max_{\sigma_{\mathcal{B}}\in S_{\mathcal{B}}}
  \sum_{A\in\mathcal{B}} g^+_\lambda(A,\sigma_{\mathcal{B}}(A)).
\]
This is an ordinary assignment problem inside the block.

\begin{lemma}[Trace-preserving positive reduction]
If
\[
  \sum_{\mathcal{B}} M_{\mathcal{B}}(\lambda)
  \geq
  m\sum_{i\in P}\lambda_i\left(\frac12-p_i\right)
  +
  m\sum_{j\in J}\lambda_j\left(\frac12-p_j\right),
\]
then the weighted critical assignment inequality holds for $\lambda$.
\end{lemma}

\begin{proof}
For each $J$-trace block $\mathcal{B}$, choose a maximizing permutation
$\sigma_{\mathcal{B}}$ in the definition of $M_{\mathcal{B}}(\lambda)$.
Combining these permutations gives a permutation $\sigma$ of
$\mathcal{F}$ preserving the $J$-trace of every set.  Hence
$q_j(\sigma)=p_j$ for all $j\in J$.  The positive gain of this permutation is
exactly $\sum_{\mathcal{B}}M_{\mathcal{B}}(\lambda)$.  Therefore
\[
  \sum_i\lambda_i(q_i(\sigma)-1/2)
  =
  \frac1m\sum_{\mathcal{B}}M_{\mathcal{B}}(\lambda)
  -
  \sum_{i\in P}\lambda_i\left(\frac12-p_i\right)
  -
  \sum_{j\in J}\lambda_j\left(\frac12-p_j\right).
\]
The assumed inequality says that the right-hand side is nonnegative, which is
the weighted critical assignment inequality.
\end{proof}

A cheaper sufficient condition comes from random block permutations.  If
$p_{i|\mathcal{B}}$ denotes the conditional frequency of coordinate $i$ inside
the block $\mathcal{B}$, then averaging a uniform random permutation of each
block gives expected positive gain
\[
  \sum_{\mathcal{B}}|\mathcal{B}|
  \sum_{i\in P}\lambda_i
  p_{i|\mathcal{B}}(1-p_{i|\mathcal{B}}).
\]
Thus the trace-preserving reduction holds whenever
\[
  \frac1m
  \sum_{\mathcal{B}}|\mathcal{B}|
  \sum_{i\in P}\lambda_i
  p_{i|\mathcal{B}}(1-p_{i|\mathcal{B}})
  \geq
  \sum_{i\in P}\lambda_i\left(\frac12-p_i\right)
  +
  \sum_{j\in J}\lambda_j\left(\frac12-p_j\right).
\]

\begin{problem}[Trace-block boost problem]
Determine when a remaining-cone weight vector satisfies the trace-preserving
positive reduction.  When it fails, quantify how much negative-coordinate
``leakage'' is needed to compensate for the missing trace-preserving boost.
\end{problem}

The trace-preserving positive reduction is not the final theorem.  It is a
useful sufficient condition, but it can fail even while the full assignment
inequality succeeds.

\begin{example}[Controlled leakage is sometimes necessary]
There is a reduced union-closed family on six coordinates with
$|\mathcal{F}|=27$ and two maximal critical coordinates $i,j$ such that
\[
  p_i=\frac8{27},
  \qquad
  p_j=\frac{11}{27}.
\]
Take weights
\[
  \lambda_i=1,
  \qquad
  \lambda_j=-\frac1{25}.
\]
This weight vector lies in the remaining cone:
\[
  \lambda_i\left(\frac12-p_i\right)
  +
  \lambda_j\left(\frac12-p_j\right)
  =
  \frac15>0,
\]
and
\[
  \lambda_i\left(2p_i-p_i^2-\frac12\right)
  +
  \lambda_j\left(2p_j-p_j^2-\frac12\right)
  =
  -\frac{7}{6075}<0.
\]
If we preserve the $j$-trace, the best boost gives
\[
  q_i=\frac{13}{27},
  \qquad
  q_j=\frac{11}{27},
\]
and the weighted slack is negative:
\[
  \lambda_i(q_i-1/2)+\lambda_j(q_j-1/2)
  =
  -\frac{2}{135}<0.
\]
However the unrestricted assignment can choose
\[
  q_i=\frac{16}{27},
  \qquad
  q_j=\frac{14}{27},
\]
giving positive weighted slack
\[
  \lambda_i(q_i-1/2)+\lambda_j(q_j-1/2)
  =
  \frac{62}{675}>0.
\]
Thus the proof cannot insist on preserving every negative coordinate.  It
must allow controlled leakage in the negative coordinates when the positive
boost gained is sufficiently larger.
\end{example}

\begin{problem}[Controlled-leakage assignment]
For a remaining-cone weight vector, prove that either the trace-preserving
positive reduction succeeds, or there is an assignment whose additional
positive gain exceeds the weighted cost of increasing the negative
coordinates.  The two-coordinate case should be attacked first, since the
example above already shows the necessary phenomenon.
\end{problem}

\subsection{Trace-Count Reduction}

The controlled-leakage problem has an important simplification: for a fixed
coordinate set, the OR-marginal polytope depends only on the counts of
coordinate traces.

Let $I$ be a finite set of coordinates.  For a trace
$u\in\{0,1\}^{I}$, write
\[
  n_u=|\{A\in\mathcal{F}: A|_I=u\}|.
\]
For a uniform coupling of two members of $\mathcal{F}$, let
$T_{u,v}$ be the amount of mass transported from trace $u$ to trace $v$,
scaled by $m=|\mathcal{F}|$.  Then $T$ is a nonnegative matrix satisfying
\[
  \sum_v T_{u,v}=n_u,
  \qquad
  \sum_u T_{u,v}=n_v.
\]
Conversely, every such transportation matrix is realized by a convex
combination of uniform couplings of $\mathcal{F}$.

\begin{proposition}[Trace-count reduction]
For a fixed coordinate set $I$, the set of possible OR-marginal vectors
$q_I(\pi)$ over all uniform self-couplings $\pi$ is exactly the image of the
transportation polytope
\[
  \{T_{u,v}\geq0:
    \sum_vT_{u,v}=n_u,\ \sum_uT_{u,v}=n_v\}
\]
under the linear map
\[
  q_i(T)
  =
  \frac1m
  \sum_{u,v:\ (u\vee v)_i=1} T_{u,v}.
\]
\end{proposition}

\begin{proof}
Any coupling of $\mathcal{F}$ pushes forward to such a transportation matrix
on $I$-traces, and the displayed formula for $q_i$ is immediate.  Conversely,
given $T$, distribute the mass $T_{u,v}$ arbitrarily among the complete
bipartite pairs between the sets with trace $u$ and the sets with trace $v$.
This gives a coupling of the uniform measure on $\mathcal{F}$ with itself.
\end{proof}

Thus the two-coordinate controlled-leakage problem is a statement about a
$4\times4$ transportation polytope, independent of union-closure except
through which trace-count vectors can occur in reduced maximal critical
families.

\begin{proposition}[Two-coordinate critical centering]
Let $I=\{i,j\}$.  If
\[
  \frac14\leq p_i<\frac12,
  \qquad
  \frac14\leq p_j<\frac12,
\]
then the trace-count transportation polytope contains a coupling with
\[
  q_i=q_j=\frac12.
\]
Equivalently, every two-coordinate weighted critical assignment inequality
holds, with no union-closure hypothesis needed beyond the realized trace
counts.
\end{proposition}

\begin{proof}
Normalize the trace counts by $m=|\mathcal{F}|$, and write
\[
  a=\Prb[00],\qquad b=\Prb[10],\qquad
  c=\Prb[01],\qquad d=\Prb[11].
\]
Thus
\[
  p_i=b+d,\qquad p_j=c+d.
\]
Start with the diagonal self-coupling.  If we move an amount $\varepsilon$
symmetrically from the two diagonal atoms $u,u$ and $v,v$ to the off-diagonal
atoms $u,v$ and $v,u$, then the two marginals are unchanged.  This swap
increases $q_i$ by $\varepsilon$ exactly when the traces $u$ and $v$ differ in
coordinate $i$, and similarly for $j$.

It is therefore enough to choose nonnegative edge weights on the complete
graph on the four traces, with incident weight at a trace at most its mass,
so that the total weight crossing the $i$-cut is
\[
  D_i=\frac12-p_i=\frac12-b-d
\]
and the total weight crossing the $j$-cut is
\[
  D_j=\frac12-p_j=\frac12-c-d.
\]
The assumptions give $0\leq D_i,D_j\leq 1/4$.

Assume first that $c\geq b$; the other case follows by interchanging
$i$ and $j$.  Then
\[
  D_i-D_j=c-b=:E.
\]
We first spend this excess $i$-demand on edges which cross only the $i$-cut.
Put
\[
  z=\min(d,E)
  \quad\text{on the edge }01\text{--}11,
  \qquad
  u=E-z
  \quad\text{on the edge }00\text{--}10.
\]
The choice of $z$ is plainly feasible.  If $u>0$, then $E>d$, so
$u=c-b-d$.  The inequality $u\leq b$ follows from
\[
  c\leq \frac12-d\leq 2b+d,
\]
where the last inequality is $b+d\geq 1/4$.  Also $u\leq a$ follows from
$c<1/2$.  Thus the excess step is feasible.  After it, the two remaining
demands are both equal to $D_j$.

There are two cases.  If $E\leq d$, then the residual masses at the traces
$00,10,01,11$ are
\[
  a,\quad b,\quad b,\quad d-E.
\]
Use $\min(b,D_j)$ on the edge $10$--$01$.  If this does not exhaust the
remaining demand, the leftover amount is
\[
  D_j-b=\frac12-b-c-d=a-\frac12.
\]
This is at most $a$, and it is at most $d-E=d-c+b$ because
\[
  \frac12-b-c-d\leq d-c+b
  \quad\Longleftrightarrow\quad
  \frac12\leq 2(b+d).
\]
The last inequality is again $p_i\geq 1/4$.  Hence the leftover can be put on
the edge $00$--$11$.

If $E>d$, then after the excess step the trace $11$ has no residual mass, and
we put all of the remaining demand $D_j$ on the edge $10$--$01$.  This fits at
trace $10$ because
\[
  D_j\leq 2b+d-c
  \quad\Longleftrightarrow\quad
  \frac12\leq 2(b+d),
\]
and it fits at trace $01$ because
\[
  D_j\leq c-d
  \quad\Longleftrightarrow\quad
  c\geq \frac14.
\]
The latter follows from $E>d$, since then $c>b+d=p_i\geq 1/4$.

In all cases we have built feasible symmetric swaps whose $i$- and $j$-cut
weights are exactly $D_i$ and $D_j$.  Applying these swaps to the diagonal
coupling gives a self-coupling with
\[
  q_i=p_i+D_i=\frac12,
  \qquad
  q_j=p_j+D_j=\frac12.
\]
Scaling back by $m$ gives the required point in the trace-count
transportation polytope.
\end{proof}

The script still contains an exhaustive sanity check for this proposition on
all two-coordinate trace count vectors with total size at most $16$.  The
check is deliberately trace-count based: it enumerates $4\times4$
transportation matrices rather than permutations of the full family.  This
makes the controlled-leakage phenomenon visible as a small transportation
fact.

\subsection{The First Higher-Dimensional Test}

The proof just given has a useful general form.  For a coordinate set $I$,
let $\mu_u$ be the normalized mass of trace $u\in\{0,1\}^I$, and put
\[
  D_i=\frac12-p_i.
\]
Starting from the diagonal self-coupling, a symmetric swap of mass
$\varepsilon$ between two traces $u$ and $v$ preserves both marginals and
increases $q_i$ by $\varepsilon$ exactly when $u_i\neq v_i$.  Therefore the
following finite certificate is enough to center all coordinates in $I$: find
edge weights $\varepsilon_{u,v}\geq0$ on unordered pairs of traces such that
\[
  \sum_{v\neq u}\varepsilon_{u,v}\leq \mu_u
  \qquad\text{for every trace }u
\]
and
\[
  \sum_{u_i\neq v_i}\varepsilon_{u,v}=D_i
  \qquad\text{for every }i\in I.
\]
The two-coordinate proposition proves that this capacitated cut-demand problem
is always feasible for $|I|=2$ when both coordinates are critical.

The naive higher-dimensional version is already false if one ignores
union-closure.
For three coordinates, take normalized trace counts with total mass $7$:
\[
  n_{100}=n_{010}=n_{001}=2,\qquad
  n_{011}=1,
\]
and all other trace counts zero.  Then
\[
  (p_1,p_2,p_3)=\left(\frac27,\frac37,\frac37\right),
\]
so all three coordinates are critical, but an exact enumeration of the
$8\times8$ trace transportation polytope finds no coupling with
$q_1=q_2=q_3=1/2$.  This is not a legal trace support for a union-closed
family: for instance $100\vee010=110$ is missing.  Thus the third coordinate
is precisely where the join-closure of the trace support has to become part
of the proof, not just the one-coordinate marginals.

A small exhaustive check first seemed to support this diagnosis: among the
$147$ three-coordinate critical trace count vectors of total size at most $9$
whose positive support is union-closed, the script found no miss.  But the
next layer produces a genuine obstruction.

\begin{proposition}[A parametric three-coordinate trace obstruction]
Let $k\geq2$.  For three coordinates $x,y,z$, take trace counts
\[
  n_{000}=4,\qquad
  n_{110}=k,\qquad
  n_{001}=k,\qquad
  n_{011}=2,\qquad
  n_{111}=1,
\]
and all other trace counts zero.  The positive trace support is union-closed,
and
\[
  (p_x,p_y,p_z)
  =
  \left(
    \frac{k+1}{2k+7},
    \frac{k+3}{2k+7},
    \frac{k+3}{2k+7}
  \right),
\]
so all three coordinates are critical.  Nevertheless the trace-count
transportation polytope does not contain
$(1/2,1/2,1/2)$.
\end{proposition}

\begin{proof}
Use the signed functional
\[
  \varphi(q)=q_x-q_y-q_z.
\]
At the desired center, $\varphi=-1/2$.  On the support
\[
  S=\{000,110,001,011,111\},
\]
define potentials
\[
  \tau(000)=0,\quad
  \tau(110)=0,\quad
  \tau(001)=-\frac12,\quad
  \tau(011)=-1,\quad
  \tau(111)=0.
\]
A direct check of the $5\times5$ table gives
\[
  \varphi(u\vee v)\leq \tau(u)+\tau(v)
  \qquad\text{for all }u,v\in S.
\]
Therefore every self-coupling of the trace distribution satisfies
\[
  \E[\varphi(U\vee V)]
  \leq
  2\E[\tau(U)]
  =
  2\cdot\frac{k(-1/2)+2(-1)}{2k+7}
  =
  -\frac{k+4}{2k+7}
  <
  -\frac12.
\]
Thus the center is separated by the assignment-dual direction
$(1,-1,-1)$.
\end{proof}

This obstruction is not merely an artificial non-realizable count vector.  It
is the projection, onto coordinates $(1,2,4)$, of the union-closed family on
five coordinates
\[
  \{\,
  0,1,7,8,9,15,24,25,28,29,31
  \,\},
\]
where sets are written as bit masks.  But that lift has hidden Frankl
witnesses: the two auxiliary coordinates have frequencies $7/11$ and $8/11$.
So the failed trace-count theorem is not a dead end; it is a sharper version
of the entropy--Hall dichotomy.  The three-coordinate dual obstruction can
exist at the projected trace level, but in this lift the obstruction is paid
for by heavy coordinates inside the fibers.

In fact this is forced for every lift of this obstruction profile.

\begin{lemma}[Bottom-fiber propagation]
Let $\mathcal{F}$ be a union-closed family, projected onto a coordinate set
$I$.  Let $A$ be the fiber over the zero trace $0\in\{0,1\}^I$, viewed as a
family of auxiliary sets in the remaining coordinates.  Suppose an auxiliary
coordinate $h$ appears in $r$ members of $A$.

If $u$ is any projected trace such that the fiber $F_u$ has a maximum
auxiliary set $M_u$, then $h\in M_u$.  Consequently, if $U$ is a collection
of such traces with nonempty fibers, then $h$ appears in at least
\[
  r+|U|
\]
members of the original family, provided $U$ does not include the zero trace.
\end{lemma}

\begin{proof}
Choose a member $a\in A$ containing $h$.  Since the projected trace of $a$ is
zero, the union of $a$ with any member of $F_u$ still has projected trace
$u$.  In particular, $a\cup M_u$ belongs to $F_u$.  But $M_u$ is maximum in
that fiber and $M_u\subseteq a\cup M_u$, so $a\cup M_u=M_u$.  Hence
$h\in M_u$.
\end{proof}

\begin{lemma}[Four-set local Frankl]
Every four-element union-closed family of auxiliary sets has an element which
appears in at least two members.
\end{lemma}

\begin{proof}
If every member were empty there would not be four distinct sets.  Suppose, to
the contrary, that every auxiliary element appears in at most one member.
Then any two nonempty members are disjoint.  But the union of two distinct
nonempty members contains every element of both, so each of those elements
appears again in the union member.  This is impossible.  Therefore there
cannot be two distinct nonempty members; with the empty set, this gives at
most two distinct members total, again impossible.  Hence some element appears
in at least two members.
\end{proof}

\begin{lemma}[Small local Frankl bounds]
The following local bounds hold for finite union-closed families of auxiliary
sets:
\[
\begin{array}{c|ccccc}
|G| & 4 & 5 & 6 & 7 & 8\\ \hline
\max_h |\{g\in G:h\in g\}| & \geq2 & \geq3 & \geq4 & \geq4 & \geq4.
\end{array}
\]
The six-set bound is deliberately stronger than the ordinary half bound.
\end{lemma}

\begin{proof}
For $|G|=4$, the preceding four-set local Frankl lemma proves the claim.  For
sizes $5$, $6$, and $7$ we use a finite incidence-pattern reduction, which is
also the intended Lean proof route.  The size-$8$ line is the ordinary
small-family half bound and is included as a separate small-case input; the
naive incidence enumeration for it is substantially larger.

Assume $G$ has $m$ members and no element reaches the claimed threshold $t$.
Let $M$ be the maximum member of $G$, and label its row by $m-1$.  Every
active auxiliary coordinate then has an incidence column containing row
$m-1$ and at most $t-2$ of the other rows.  There are only
\[
  \sum_{j=0}^{t-2}\binom{m-1}{j}
\]
such columns.  The script enumerates every subset of these columns, keeps the
ones that separate the $m$ rows, and checks whether the resulting row set is
closed under coordinatewise OR.  For
\[
  (m,t)=(5,3),(6,4),(7,4)
\]
there is no such row set.  This is a finite certificate, printed by
\texttt{--local-frankl-checks}.
\end{proof}

\begin{lemma}[Horn propagation certificate]
Fix an auxiliary coordinate $h$ obtained from the zero fiber.  For each trace
u, let $y_u$ be the number of members in the $u$-fiber containing $h$, and
let $n_u$ be the fiber size.  Then the following constraints are sound:
\[
  y_0\geq r_0,
  \qquad
  y_u\geq1 \quad (u\neq0,\ n_u>0),
\]
where $r_0$ is any valid local lower bound in the zero fiber, and
\[
  y_{u\vee v}=n_{u\vee v}
  \quad\Longrightarrow\quad
  y_u=n_u\ \text{or}\ y_v=n_v.
\]
\end{lemma}

\begin{proof}
The first inequality is the chosen local zero-fiber bound.  The second is
bottom-fiber propagation: each nonzero fiber maximum contains $h$.

For the Horn implication, suppose all members of the $(u\vee v)$-fiber
contain $h$.  If some member of the $u$-fiber omitted $h$ and some member of
the $v$-fiber also omitted $h$, then their union would omit $h$ but would lie
in the $(u\vee v)$-fiber, a contradiction.  Hence one of the two input
fibers must consist entirely of members containing $h$.
\end{proof}

\begin{proposition}[Unique-top symbolic propagation]
Let $\mathcal{F}$ be a union-closed lift of a three-coordinate trace profile
whose nonzero support consists of four traces $B,C,D,T$ satisfying
\[
  B\vee C=T,\qquad B\vee D=T,\qquad C\vee D=D,
\]
and suppose the $T$-fiber has exactly one member.  Let the fiber sizes be
\[
  |A|=z,\qquad |B|=b,\qquad |C|=c,\qquad |D|=d,\qquad |T|=1,
\]
where $A$ is the zero fiber.  If an auxiliary coordinate $h$ appears in
$r$ members of $A$, then $h$ appears in at least
\[
  r+\min(b+3,\;c+d+2)
\]
members of $\mathcal{F}$.
\end{proposition}

\begin{proof}
By bottom-fiber propagation, $h$ appears in the maximum member of each
nonzero fiber $B,C,D,T$.  If every member of $B$ contains $h$, then together
with one occurrence in each of $C,D,T$ we get at least
\[
  r+b+3
\]
occurrences.

Otherwise choose $b_0\in B$ with $h\notin b_0$.  Since $B\vee C=T$ and the
$T$-fiber is a singleton, $b_0\cup c$ is the unique member of $T$ for every
$c\in C$.  The unique member of $T$ contains $h$ by bottom-fiber propagation,
so every member of $C$ contains $h$.  The same argument using
$B\vee D=T$ shows that every member of $D$ contains $h$.  Together with one
member of $B$ and the unique member of $T$, this gives at least
\[
  r+c+d+2
\]
occurrences.  Taking the worse of the two cases proves the claim.
\end{proof}

\begin{proposition}[The parametric obstruction forces a hidden witness]
Let $k\geq2$.
Let $\mathcal{F}$ be a union-closed family whose projection onto three
coordinates has trace counts
\[
  n_{000}=4,\qquad
  n_{110}=k,\qquad
  n_{001}=k,\qquad
  n_{011}=2,\qquad
  n_{111}=1,
\]
with all other trace counts zero.  If all members with the same trace are
viewed as auxiliary sets in the remaining coordinates, then some auxiliary
coordinate belongs to at least $k+5$ of the $2k+7$ members of
$\mathcal{F}$.  In particular, it is a Frankl witness.
\end{proposition}

\begin{proof}
Let $A$ be the $000$ fiber.  Since $000\vee000=000$, the four auxiliary sets
in $A$ form a union-closed family.  By the four-set local Frankl lemma, this
fiber has an element $h$ which appears in at least two members of $A$.
Apply the unique-top symbolic propagation proposition with
\[
  r=2,\qquad b=k,\qquad c=k,\qquad d=2.
\]
It gives at least
\[
  2+\min(k+3,k+4)=k+5
\]
occurrences, while
\[
  \frac{|\mathcal{F}|}{2}=\frac{2k+7}{2}<k+5.
\]
\end{proof}

\begin{corollary}[The first sign obstructions are fiber-paid]
For three-coordinate critical trace counts of total size at most $16$, the
script's propagation scan
\[
  \texttt{--trace-three-propagation-scan-size 16}
\]
finds that every $\{-1,0,1\}$-separated class is paid for by the hidden
element propagation certificate above.
\end{corollary}

\begin{proof}
This is a finite computation.  For each separated trace-count vector, the
script minimizes $\sum_u y_u$ subject to the small local zero-fiber bounds,
bottom-fiber propagation, and the Horn implications above.  In every case of
total size at most $16$, the minimum is at least half the total family size.
\end{proof}

At total size $17$, the same scan reports the first unpaid classes.  With the
size-$8$ local bound included, the four canonical unpaid classes are
\[
  (6,0,0,3,3,0,3,2),
\]
\[
  (8,0,1,0,1,0,1,6),\quad
  (8,0,1,0,1,0,2,5),\quad
  (8,0,1,0,2,0,1,5).
\]
Each has propagation lower bound $8$, exactly one occurrence short of half.
Starting the hidden element in a nonzero fiber does not improve this bound.
This was the next obstruction layer for the one-coordinate propagation
certificate.  It is closed by a slightly more fiber-sensitive shadow
certificate.

A tempting strengthening would be: if the zero fiber has six members and
$h$ appears in four of them, then any three-member fiber preserved by the
zero fiber must contain $h$ in at least two members.  This is false in
isolation.  For example, with auxiliary bit masks
\[
  A=\{0,1,3,7,11,15\},
  \qquad
  X=\{12,13,15\},
\]
the family $A$ is union-closed, the coordinate corresponding to bit $1$
appears in four members of $A$, the three-set family $X$ is union-closed and
is preserved by union with $A$, but that bit appears in only one member of
$X$.  Hence the missing size-$17$ occurrence cannot be forced by a
single-fiber action lemma alone; it must use the simultaneous constraints
among the several nonzero fibers, or else a different transport idea.

\begin{definition}[Zero-fiber shadow]
Let $A$ be the zero fiber and let $J$ be a set of auxiliary coordinates that
separates the members of $A$, after quotienting clone coordinates.  For any
other fiber $F_u$, its $J$-shadow is the multiset of restrictions
\[
  \{\,x|_J:x\in F_u\,\}.
\]
We allow repeated shadows, because members which agree on $J$ may still be
separated by auxiliary coordinates outside $J$.
\end{definition}

The point of this relaxation is that it can only make counterexamples easier
to build.  The zero fiber acts on every nonzero fiber by union, so every
$J$-shadow must be preserved by union with the $J$-shadow of $A$.  Also, if
$u\vee v=w$, then every union of a member of the $u$-shadow with a member of
the $v$-shadow must lie in the support of the $w$-shadow.  These are only
necessary conditions for a genuine lift, but they remember simultaneous
constraints between the nonzero fibers.

\begin{proposition}[The size-$17$ frontier is shadow-paid]
For each of the four size-$17$ unpaid trace-count classes above, every
union-closed lift has an auxiliary coordinate, already visible in the zero
fiber, which appears in at least $9$ of the $17$ members.
\end{proposition}

\begin{proof}
First suppose the zero fiber contains an auxiliary coordinate appearing in at
least five zero-fiber members.  In all four frontier classes there are four
nonzero trace fibers.  Bottom-fiber propagation puts that coordinate in the
maximum of each nonzero fiber, so it appears at least $5+4=9$ times.

It remains to handle the equality cases where every zero-fiber coordinate has
frequency at most four.  For a six-member zero fiber, the reduced
union-closed shadows with maximum frequency at most four are finite; the
script enumerates the $13$ canonical shadows.  For an eight-member zero
fiber, the relevant equality shadow, after quotienting clone coordinates, is
the Boolean three-cube.  For each zero shadow the script enumerates every
multiset shadow of the nonzero fibers which is preserved by the zero fiber
and satisfies all trace-join support constraints.

The command
\[
  \texttt{--trace-three-frontier-shadow}
\]
reports the following minimum, over this relaxed shadow model, of the largest
frequency among zero-fiber coordinates:
\[
\begin{array}{c|c}
\text{trace counts} & \text{forced zero-coordinate frequency}\\ \hline
(6,0,0,3,3,0,3,2) & 12\\
(8,0,1,0,1,0,1,6) & 11\\
(8,0,1,0,1,0,2,5) & 11\\
(8,0,1,0,2,0,1,5) & 11.
\end{array}
\]
All four lower bounds exceed $9$.  Since the shadow model permits repeated
shadows and forgets all auxiliary coordinates outside the zero fiber, every
genuine lift is a special case of the relaxation.  Therefore each frontier
class has a hidden Frankl witness.
\end{proof}

The integrated scan
\[
  \texttt{--trace-three-propagation-scan-size N}
\]
now runs this as a second layer after the Horn propagation certificate and
prints the first remaining frontier, if any.  At size $17$ it reports $48$
propagation-paid separated counts and four shadow-paid canonical classes.  A
run through size $18$ finds no remaining frontier.

The symbolic task is now organized by support templates rather than by total
size.  The command
\[
  \texttt{--trace-three-template-scan-size 18}
\]
finds five canonical three-coordinate support templates among all separated
critical trace counts through size $18$:
\[
\begin{gathered}
  \{0,1,2,3,7\},\quad
  \{0,1,3,6,7\},\quad
  \{0,3,5,6,7\},\\
  \{0,1,2,3,5,7\},\quad
  \{0,1,3,5,6,7\}.
\end{gathered}
\]
The unique-top symbolic propagation proposition covers one recurring
five-trace template whenever the top fiber is a singleton.  The next
symbolic objective is to prove comparable parametric certificates for the
remaining support templates, with finite shadow checks reserved only for
their small equality cases.

The refined target is therefore not a pure trace-count theorem.  It is a
fiber-sensitive statement: if a projected three-coordinate obstruction of
this kind occurs in an actual union-closed family, then either controlled
transport using the internal fibers restores centering, or the same dual
certificate can be lifted into a genuine Frankl witness outside the projected
coordinates.  The command \texttt{--trace-three-obstruction} prints this
trace obstruction, its separating direction, and the explicit lift above.

\subsection{Lean-Ready Boundary}

The part of the argument that now looks stable enough to formalize is not the
whole Frankl strategy, but the first obstruction layer.  It breaks into a
small reusable API:
\begin{enumerate}
  \item A finite nonempty union-closed family has a maximum member, namely the
  union of all its members.
  \item Bottom-fiber propagation: if a zero-trace auxiliary coordinate appears
  in a bottom fiber, then it appears in the maximum member of every nonzero
  fiber.
  \item Small local Frankl bounds for zero fibers of sizes $4$ through $8$.
  \item The Horn propagation certificate for fibers whose joins land in
  all-$h$ fibers.
  \item The unique-top symbolic propagation proposition, which turns a whole
  parametric trace family into one inequality.
  \item The parametric three-coordinate trace obstruction is separated by the
  explicit assignment-dual potential above.
  \item Every union-closed lift of that parametric obstruction has a hidden
  Frankl witness.
  \item The four size-$17$ frontier classes are paid by the zero-fiber shadow
  certificate.
\end{enumerate}
These are local, finite, and do not depend on any unproved global form of
Frankl's conjecture.  They are therefore good candidates for the next Lean
pass before attempting to formalize the broader entropy--transport program.

\subsection{Coordinate-Order Faces}

Some tight dual directions are structural and should be stripped away before
looking for a real obstruction.  Define a preorder on coordinates by
\[
  i\preceq j
  \qquad\Longleftrightarrow\qquad
  \forall A\in\mathcal{F},\ i\in A\Rightarrow j\in A.
\]
Equivalently, $\chi_i\leq \chi_j$ pointwise.  If $i\preceq j$, then every
coupled OR marginal satisfies
\[
  q_i(\pi)\leq q_j(\pi).
\]
Indeed, $i\in A\cup B$ implies $i\in A$ or $i\in B$, hence $j\in A$ or
$j\in B$, so $j\in A\cup B$.

Thus the hyperplane $q_i=q_j$ can support the centering polytope at
$(1/2,\dots,1/2)$ without representing a failure.  The corresponding signed
dual direction is $\lambda_i=1$, $\lambda_j=-1$.  Clone coordinates give the
special case $i\preceq j$ and $j\preceq i$, where $q_i=q_j$ for all couplings.

This suggests the following normalization before attacking the assignment
dual:
\[
  \text{quotient coordinate clones, then keep only maximal coordinates under }
  \preceq.
\]
For Frankl this is harmless as a search principle: if $i\preceq j$, then
$p_i\leq p_j$, so the dominated coordinate $i$ is never a better witness than
$j$.  In a putative counterexample all maximal coordinates are still light.

\begin{problem}[Face-stripped assignment inequality]
Prove the weighted critical assignment inequality after quotienting clone
coordinates and restricting to maximal critical coordinates.  Then solve the
extension problem for dominated coordinates with compatible targets
$q_i\leq q_j$ when $i\preceq j$.  Exact centering of every dominated
coordinate may be more than the entropy argument needs; a target inside
$p_i<q_i<1-p_i$ would still give coordinate entropy gain.
\end{problem}

\section{The First Novel Program}

\begin{problem}[Optimal coupled OR]
Among all couplings $\pi\in\Pi(\mathcal{F})$, maximize
\[
  \Phi(\pi)=\Ent(X\vee Y)-\Ent(X).
\]
Understand the dual certificate when the maximum is nonpositive.
\end{problem}

This problem is nonlinear because of entropy, but it has a transport flavor:
we are choosing mass on $\mathcal{F}\times\mathcal{F}$ with fixed marginals,
and the objective depends only on the OR fibers
\[
  m^{-1}(C)=\{(A,B): A\vee B=C\}.
\]
The failure of entropy gain should therefore be witnessed by a potential on
OR fibers and coordinates.

\begin{heuristic}[Entropy failure equals collision structure]
If no coupling gives positive entropy advantage, then the OR map
\[
  m:\mathcal{F}\times\mathcal{F}\to\mathcal{F},
  \qquad
  m(A,B)=A\vee B,
\]
has enough forced collision/dependence structure to produce a combinatorial
certificate.  That certificate should be interpretable as a Hall obstruction
to naive pairing.
\end{heuristic}

\section{Hall Certificates for Pairing Failure}

Fix a coordinate $i$.  Let
\[
  L_i = \{A\in\mathcal{F}: i\notin A\},
  \qquad
  R_i = \{B\in\mathcal{F}: i\in B\}.
\]
Frankl for coordinate $i$ is exactly
\[
  |L_i| \leq |R_i|.
\]

Build a bipartite graph $G_i$ from $L_i$ to $R_i$ by declaring that
$A\in L_i$ is adjacent to $B\in R_i$ if
\[
  B=A\cup C
  \quad \text{for some } C\in\mathcal{F} \text{ with } i\in C.
\]
Union-closure guarantees $B\in\mathcal{F}$ whenever $C$ is chosen.

\begin{lemma}[Matching certificate, desired]
If $G_i$ has a matching saturating $L_i$, then $i$ is a Frankl element.
\end{lemma}

This is immediate from Hall's theorem.  The point is what happens when no such
matching exists.

\begin{definition}[Hall defect]
For $U\subseteq L_i$, define
\[
  \delta_i(U)=|U|-|N_i(U)|.
\]
A positive Hall defect is a certificate that the obvious union-pairing strategy
for coordinate $i$ cannot work.
\end{definition}

In a counterexample, every coordinate has $|L_i|>|R_i|$, so every coordinate
has global counting pressure against Frankl.  But the graphs $G_i$ are not
independent: all are induced by the same join operation on $\mathcal{F}$.

\begin{problem}[Incompatibility of Hall defects]
Show that a finite union-closed family cannot support coherent positive Hall
defects for every active coordinate, unless some coordinate has
$p_i\geq 1/2$.
\end{problem}

\section{The Entropy--Hall Synthesis}

The central proposal is a dichotomy.

\begin{conjecture}[Entropy--Hall dichotomy]
Let $\mathcal{F}$ be a finite union-closed family with $0<p_i<1/2$ for every
active coordinate.  Then at least one of the following holds:
\begin{enumerate}[label=(\alph*)]
  \item There exists a coupling $\pi\in\Pi(\mathcal{F})$ with
  $\Phi(\pi)>0$.
  \item The dual obstruction to $\Phi(\pi)>0$ yields Hall-defect data
  $\{\delta_i(U_i)>0\}_i$ that is incompatible with union-closure.
\end{enumerate}
\end{conjecture}

The first alternative gives an entropy contradiction.  The second alternative
is meant to convert the entropy method's slack into a purely finite
combinatorial contradiction.

\section{Encountering the Barriers One by One}

\subsection{Barrier: iid OR stops around \texorpdfstring{$(3-\sqrt{5})/2$}{(3-sqrt(5))/2}}

Path through: replace iid OR by a positively correlated coupling.  The target
is not arbitrary dependence, but dependence tuned so that
\[
  p_i < q_i(\pi) < 1-p_i
\]
for all coordinates.  This preserves positive coordinate entropy gain while
still moving away from the original distribution.

Concrete test:
\[
  \text{Can every putative counterexample admit such a coupling?}
\]
If not, what is the transport dual obstruction?

\subsection{Barrier: entropy gains are eaten by dependence}

Path through: measure the loss explicitly as
\[
  \TC(Z)-\TC(X).
\]
This loss is not noise; it records collision structure of the OR map.  Try to
show that large dependence loss implies a strong Hall-defect pattern.

Concrete target:
\[
  \TC(Z)-\TC(X)
  \leq
  \text{a weighted sum of local collision defects}.
\]
If the right side is too large, extract a coordinate with an actual matching
or fractional matching.

\subsection{Barrier: average-size bounds lose dimension}

Path through: introduce an effective dimension for the separating coordinate
system
\[
  \{\chi_i:\mathcal{F}\to\mathbf{2}\}.
\]
Possible choices:
\begin{enumerate}[label=(\roman*)]
  \item rank of the incidence matrix over $\mathbb{Q}$;
  \item entropy dimension, $\sum_i h(p_i)$ normalized by $\log|\mathcal{F}|$;
  \item number of join-irreducible coordinates after quotienting clones;
  \item VC or shatter dimension of the coordinate traces.
\end{enumerate}
Then prove a trichotomy:
\[
  \text{dense} \quad \vee \quad
  \text{low effective dimension} \quad \vee \quad
  \text{entropy--transport contradiction}.
\]

\subsection{Barrier: dense cases are not the hard cases}

Path through: use known density theorems as a boundary condition, but focus
the new argument on sparse families.  Sparse families have larger OR-collision
fibers, which is exactly the regime where the entropy/Hall synthesis should
have more structure to exploit.

\subsection{Barrier: equivalent formulations move the problem}

Path through: choose the formulation by the obstruction it exposes.
\[
\begin{array}{ccl}
\text{set families} &\leadsto& \text{frequency and pairing},\\
\text{join-semilattices} &\leadsto& \text{coordinate homomorphisms},\\
\text{graphs} &\leadsto& \text{maximal stable set counts},\\
\text{closure systems} &\leadsto& \text{Galois duality and irreducibles}.
\end{array}
\]
The proposed note lives in the join-semilattice/probability hybrid because it
keeps both the OR operation and the counting measure visible.

\section{Concrete Lemmas to Try}

\begin{lemma}[Coupled OR window]
Assume $\max_i p_i<1/2$.  Either there exists a coupling
$\pi\in\Pi(\mathcal{F})$ such that
\[
  p_i < q_i(\pi) < 1-p_i
  \quad\text{for all active } i,
\]
or the transport dual yields a nonzero signed measure on coordinates that
separates the family from this window.
\end{lemma}

\begin{lemma}[Window plus bounded correlation loss]
If a coupling $\pi$ satisfies the OR window and
\[
  \TC(X\vee Y)-\TC(X)
  <
  \sum_i \bigl(h(q_i(\pi))-h(p_i)\bigr),
\]
then $\mathcal{F}$ satisfies Frankl.
\end{lemma}

This lemma is tautological, but useful: it isolates the only remaining
obstacle as correlation growth.

\begin{lemma}[Correlation loss gives Hall defect]
For an entropy-maximizing coupling $\pi$, large positive
$\TC(X\vee Y)-\TC(X)$ forces a positive weighted combination of Hall defects
in the graphs $G_i$.
\end{lemma}

This is the first genuinely speculative lemma.  It asks us to turn dependence
created by OR-collisions into a finite matching obstruction.

\begin{lemma}[Hall defects uncross]
For each coordinate $i$, among all positive Hall defects choose an inclusion
minimal $U_i\subseteq L_i$.  The family of minimal defects can be chosen to
satisfy an uncrossing property compatible with joins:
\[
  U_i \vee U_j
  =
  \{A\vee B:A\in U_i,\ B\in U_j\}
\]
has controlled neighborhood growth.
\end{lemma}

\begin{lemma}[No coherent defect web]
No reduced finite join-semilattice with separating coordinates and
$p_i<1/2$ for all $i$ admits a coherent system of uncrossed positive Hall
defects for every coordinate.
\end{lemma}

If the last two lemmas are true in an appropriately sharpened form, they would
complete the entropy--Hall dichotomy.

\section{A Minimal-Counterexample Normalization}

Suppose $\mathcal{F}$ is a counterexample of minimal size.  We should be able
to enforce the following normalizations:
\begin{enumerate}[label=(N\arabic*)]
  \item The ground set is exactly $\bigcup\mathcal{F}$.
  \item Distinct coordinates induce distinct homomorphisms
  $\chi_i:\mathcal{F}\to\mathbf{2}$.
  \item No singleton belongs to $\mathcal{F}$, by the already formalized
  singleton case.
  \item Every coordinate has $0<p_i<1/2$.
  \item The family is closed under deleting irrelevant duplicate coordinates.
\end{enumerate}
Each normalization should be formalized separately.  If any normalization
fails, it may already yield a smaller counterexample or a Frankl element.

\section{Formalization Roadmap}

The Lean project now has a small checked spine for the first obstruction
layer.  The following table separates kernel-checked content from the remaining
mathematical bridge lemmas.

\begin{center}
\begin{tabular}{p{0.22\linewidth}p{0.24\linewidth}p{0.44\linewidth}}
\textbf{Layer} & \textbf{Lean artifact} & \textbf{Current status} \\
\hline
Basic Frankl reductions &
\texttt{Frankl.Basic}, \texttt{Frankl.Split} &
Singleton-containing families and the split-family witness equivalence are
checked. \\
Coupled one-coordinate algebra &
\texttt{Frankl.CoupledOr}, \texttt{Frankl.Probability},
\texttt{Frankl.Entropy}, \texttt{Frankl.EntropyGain},
\texttt{Frankl.TotalCorrelation}, \texttt{Frankl.CoupledHinge} &
The Frechet interval, critical centering, finite two-Bernoulli couplings, and
the entropy gain/total-correlation accounting formats are checked. \\
AHS hinge payload &
\texttt{Frankl.AHSFormula}, \texttt{Frankl.AHSData},
\texttt{Frankl.Certificates}, \texttt{Frankl.AHSInterval} &
The formulae, rational finite-row checks, and interval-row assembly lemmas are
in Lean; the logarithmic interval enclosures and derivative identity remain
open. \\
Trace frontier certificates &
\texttt{Frankl.Trace}, \texttt{Frankl.Fiber},
\texttt{Frankl.FrontierData} &
Trace-fiber propagation and the finite size-$17$ relaxed shadow predicates are
checked; realizing the full shadow search from first principles remains open. \\
Symbolic propagation &
\texttt{Frankl.Boost}, \texttt{Frankl.UniqueTop} &
The arithmetic boost core and a five-fiber unique-top propagation theorem are
checked; connecting arbitrary projected families to this profile API remains a
formalization target. \\
Global coupled-kernel program &
\texttt{Frankl.CoupledHinge}, \texttt{Frankl.GlobalCoupling} &
Uniform global couplings, OR marginals, and the formal criterion are stated,
but existence of compatible global kernels and the total-correlation bound
remain the main frontier. \\
\end{tabular}
\end{center}

The computational artifacts are synchronized by the following lightweight
commands.

\begin{center}
\begin{tabular}{p{0.24\linewidth}p{0.34\linewidth}p{0.32\linewidth}}
\textbf{Artifact} & \textbf{Command} & \textbf{Status} \\
\hline
\texttt{ahs\_hinge\_certificate.json} &
\texttt{gilmer\_ahs\_verification.py --check-ahs-certificate
research/ahs\_hinge\_certificate.json} &
Checks the deterministic mpmath interval audit payload; not yet a
kernel-checked transcendental proof. \\
\texttt{center\_generation\_certificate.json} &
\texttt{window\_experiment.py --check-center-certificate
research/center\_generation\_certificate.json} &
Checks stored generated-feasible assignment-column payloads. \\
\texttt{reduced\_maximal\_center\_generation\_certificate.json} &
\texttt{window\_experiment.py --check-center-certificate
research/reduced\_maximal\_center\_generation\_certificate.json} &
Checks the reduced/maximal scan payload. \\
Size-$17$ frontier shadows &
\texttt{window\_experiment.py --campaign frontier-shadow} &
Regenerates the four relaxed zero-shadow frontier certificates now mirrored in
\texttt{Frankl.FrontierData}. \\
Total-correlation accounting &
\texttt{window\_experiment.py --campaign tc-exact-n4} &
Runs exact permutation TC accounting for small critical-center cases through
$n=4$, bounded to $|\mathcal{F}|\leq7$; the full size-$8$ exact pass is a
longer campaign. \\
Campaign presets &
\texttt{window\_experiment.py --list-campaigns} &
Lists reproducible computational campaigns used by the note. \\
\end{tabular}
\end{center}

The next formal targets should be low-risk statements that clarify the attack:
\begin{enumerate}
  \item Continue the one-coordinate coupled-OR algebra now started in
  \texttt{Frankl/CoupledOr.lean}: promote the Frechet interval calculation
  from raw real algebra to a finite probability statement about two
  Bernoulli coordinates with equal marginals.
  \item Define the split families $L_i$ and $R_i$ and prove
  $i$ is Frankl iff $|L_i|\leq |R_i|$.
  \item Define the bipartite graph $G_i$ and prove the matching certificate
  using a finite Hall theorem from Mathlib, or first prove the fractional
  version via finite sums.
  \item Define coordinate maps $\chi_i$ as join-homomorphisms to `Bool' or
  `Fin 2', and prove they separate points after reducing the ground set.
  \item Formalize the entropy identity
  \[
    \Ent(Z)-\Ent(X)
    =
    \sum_i(h(q_i)-h(p_i))-(\TC(Z)-\TC(X))
  \]
  in a separate finite-probability file, likely before connecting it to
  Frankl.
  \item Experiment computationally with small union-closed families, solving
  the coupled OR window as a linear feasibility problem.
\end{enumerate}

\section{What Would Count as Progress?}

We do not need the full dichotomy immediately.  Any of the following would be
meaningful progress:
\begin{enumerate}
  \item Find a natural class of sparse families where the coupled OR window
  exists and beats correlation growth.
  \item Derive an explicit dual certificate for failure of the OR window.
  \item Show that the dual certificate implies one of the known structural
  special cases.
  \item Prove that minimal counterexamples must have a strong collision
  profile under the OR map.
  \item Connect OR-collision fibers to Hall defects in one nontrivial family
  class.
\end{enumerate}

\section{First Computational Probe}

The companion script
\[
  \texttt{window\_experiment.py}
\]
tests the coupled-OR window on small union-closed families.  For exact small
families it uses the Birkhoff-von Neumann theorem: every uniform coupling of
two uniform samples from $\mathcal{F}$ is a convex combination of permutation
couplings.  The script enumerates the corresponding OR-marginal vectors
\[
  q(\sigma)_i
  =
  \frac{1}{|\mathcal{F}|}
  \left|\{A\in\mathcal{F}: i\in A\cup\sigma(A)\}\right|
\]
and solves the low-dimensional convex-hull LP
\[
  \max \gamma
  \quad\text{subject to}\quad
  p_i+\gamma\le q_i\le 1-p_i-\gamma.
\]

Two coordinate modes are useful:
\begin{enumerate}
  \item \texttt{nonheavy}: coordinates with $p_i\le 1/2$;
  \item \texttt{light}: coordinates with $p_i<1/2$.
\end{enumerate}
The first mode exposes exact boundary obstructions.  The second mode tests
whether strictly light coordinates can be pushed into the entropy window.

For all union-closed families on three ground elements, the exact
\texttt{light} scan found no obstruction: every tested light-coordinate window
had positive margin.  The smallest margin occurred at $p=3/7$, where the
solver found $q=1/2$ and $\gamma=1/14$.  For rarer coordinates the optimum is
often $q=2p$ rather than $q=1/2$, exactly as predicted by the
single-coordinate range.  In \texttt{nonheavy} mode, the smallest margins are
exactly zero and occur at the expected sharp boundary $p_i=1/2$, including the
Boolean cube itself.

For all union-closed families on four ground elements, with exact permutation
couplings for $|\mathcal{F}|\le 8$ and sampled permutation couplings for larger
families, the same pattern persisted.  In a critical-coordinate scan
($1/4\le p_i<1/2$), all 1988 exact checks were center-feasible, and the sampled
larger cases also found centers.  The hardest light-coordinate examples had
$p=7/15$ and sampled couplings reaching $q=1/2$, giving $\gamma=1/30$.

The script now also has a stronger \texttt{--scan-subsets} mode: instead of
testing only the highest-frequency eligible coordinates, it tests every
eligible coordinate subset of size at most three.  In the exhaustive
four-coordinate critical scan, this produced 4468 center checks:
\[
  2924\text{ exact-feasible},\qquad
  1544\text{ sampled-feasible},\qquad
  0\text{ center misses}.
\]
This is better evidence for the ``obstruction-critical subcollection''
phrasing of the conjecture than the earlier top-coordinate scan.  Two small
random subset scans also found no misses: 56 critical subset checks from 40
random five-coordinate generated families, and 37 critical subset checks from
30 random six-coordinate generated families.

The direct signed-weight stress test reinforces the face-stripping picture.
With \texttt{--lambda-stress}, the script samples signed vectors $\lambda$ and
solves the exact assignment problem for each one.  It now also reports the two
necessary cone quantities from the diagonal and random-permutation baselines,
so small raw slack is not mistaken for a genuine candidate obstruction.  On
the exhaustive four-coordinate critical subset scan, no sampled direction
separated the center.  Before normalization, many zero-slack directions were
explained by clone coordinates or by coordinate-order faces $i\preceq j$,
where $q_i\leq q_j$ is forced for every coupling.  After imposing both
\texttt{--reduced-only} and \texttt{--maximal-only}, the exhaustive
four-coordinate scan had 628 center checks and no misses; the hardest sampled
dual slack was
\[
  \frac1{30},
\]
coming from a single maximal coordinate with $p=7/15$.  A reduced, maximal
random six-coordinate scan produced only three maximal critical coordinate
sets in the sample, again all single-coordinate, with hardest slack
$1/38$.  A broader random seven-coordinate scan produced 31 reduced/maximal
critical checks, including a two-coordinate maximal critical case, again with
no misses and no sampled separating direction; the hardest slack was $1/70$.
This suggests that a large part of the apparent signed-dual difficulty is
caused by forced faces rather than genuine transport obstructions.

A targeted search for remaining-cone geometry found mixed two-coordinate
examples.  One reduced six-coordinate family has maximal critical coordinates
with frequencies
\[
  (p_i,p_j)=\left(\frac8{29},\frac{14}{29}\right).
\]
The best cone direction found had one positive low-critical weight on $i$ and
one negative high-critical weight on $j$.  The exact assignment still had
positive slack, and the maximizing permutation achieved
\[
  q_i=\frac{16}{29}=2p_i,
  \qquad
  q_j=\frac{14}{29}=p_j.
\]
In other words, it boosted the rare positive coordinate all the way to its
single-coordinate maximum while preserving the negative coordinate at the
diagonal value.  The protected-block condition explains this: inside the two
$j$-trace blocks, the conditional frequencies of $i$ were $2/15$ and $6/14$,
both at most $1/2$.

The stronger partial-boost lemma is needed because the protected-block
condition is not automatic.  A later search found a one-positive cone
direction with
\[
  (p_j,p_i)=\left(\frac{19}{39},\frac{17}{39}\right),
\]
where $j$ has negative weight and $i$ has positive weight.  In the
$j=1$ trace block, the conditional frequency of $i$ was $13/19>1/2$, so full
boosting while preserving $j$ was impossible.  Nevertheless the
trace-preserving boost size was
\[
  \Delta_{i|j}=\frac{10}{39},
\]
and the exact assignment achieved
\[
  q_j=\frac{19}{39}=p_j,
  \qquad
  q_i=\frac{27}{39}=p_i+\Delta_{i|j}.
\]
For the tested cone direction the required boost was only about $0.0033$,
so the partial boost had large slack.  This suggests that the right local
quantity is not whether every block has conditional frequency at most $1/2$,
but whether the aggregate trace-preserving boost $\Delta_{i|J}$ exceeds the
signed demand imposed by the negative coordinates.

A broader eight-coordinate reduced/maximal scan produced four remaining-cone
candidate directions.  All four were handled by the trace-preserving positive
assignment.  In the hardest one, the critical frequencies were
\[
  (p_i,p_j)=\left(\frac{21}{47},\frac{23}{47}\right),
\]
with positive weight on $i$ and negative weight on $j$.  The trace-preserving
assignment kept
\[
  q_j=p_j
\]
and boosted
\[
  q_i=\frac{42}{47}=2p_i.
\]
The required boost was only about $0.005$, while the available
trace-preserving boost was $21/47$.  This shows why the trace-preserving
criterion remains valuable: many remaining-cone directions have tiny diagonal
slack, but the available positive boost can still be enormous.  The later
controlled-leakage example shows that this criterion is not universal.

For random generated families on six ground elements, the sampled search again
found positive window margins on the most critical light coordinates.  Exact
centering was usually found as the sampled permutation hull became denser, but
some sampled ``not seen feasible'' cases remained.  These are not certificates
of failure; they are candidate stress tests for either a genuine dual
obstruction or a more systematic sampling/LP method.  The script can print the
first such families with \texttt{--show-misses}.

For a sampled center miss, the script extracts a separating coordinate weight
vector $\lambda$ from the sampled $q$-hull.  It then checks this separator
against the full permutation-coupling polytope by solving the exact assignment
problem
\[
  \max_{\sigma\in S_{\mathcal{F}}}
  \frac1{|\mathcal{F}|}
  \sum_{A\in\mathcal{F}}\sum_{i\in I}
  \lambda_i\,\mathbf{1}_{i\in A\cup\sigma(A)}.
\]
If this exact maximum is still below
\[
  \sum_{i\in I}\lambda_i/2,
\]
then the separator is a genuine obstruction to simultaneous centering for the
chosen coordinate set.  In the first printed six-coordinate stress cases, the
assignment check beat the sampled separator, so those misses were sampling
weakness rather than real obstructions.  This is encouraging: it suggests that
many apparent misses are caused by a thin sampled hull, while also giving a
clear route for recognizing a true obstruction if one appears.

The script also implements a column-generation version of this idea.  Starting
from sampled permutation couplings, it repeats:
\[
\begin{array}{ll}
  \text{(1)} & \text{if }(1/2,\dots,1/2)\text{ is in the current hull, stop;}\\
  \text{(2)} & \text{otherwise extract a separating }\lambda;\\
  \text{(3)} & \text{solve the exact assignment problem for }\lambda;\\
  \text{(4)} & \text{if the exact optimum still separates, certify obstruction;}\\
  \text{(5)} & \text{otherwise add the maximizing permutation as a new column.}
\end{array}
\]
On a stress run of 120 random six-coordinate generated families, using only
300 sampled permutations as the initial hull, all sampled critical-center
misses were resolved by column generation.  The resulting status counts were
\[
  3\text{ exact-feasible},\quad
  77\text{ sampled-feasible},\quad
  10\text{ generated-feasible},\quad
  0\text{ verified obstructions}.
\]
In the hardest printed case, the center was certified after adding just one
exact assignment column.  This strengthens the experimental evidence for
critical-coordinate centering and gives a concrete finite procedure for
searching for genuine obstructions.

These computations do not prove the coupled window lemma.  They do suggest a
sharper experimental conjecture:

\begin{conjecture}[Critical light-coordinate centering]
Let $\mathcal{F}$ be finite and union-closed.  For every set of light
coordinates in the critical range
\[
  I=\{i:1/4\le p_i<1/2\},
\]
or at least for every obstruction-critical subcollection of them, there exists
a uniform coupling of $X,Y$ such that
\[
  \Prb[(X\vee Y)_i=1]=\frac12
  \qquad\text{for all } i\in I.
\]
\end{conjecture}

For lighter coordinates $p_i<1/4$, the right target is not $1/2$ but some
value in the interval
\[
  p_i<q_i\le 2p_i<1/2,
\]
which still gives positive coordinate entropy gain.  If the critical
centering conjecture is true, then the iid entropy barrier is removed exactly
where iid OR fails.  The remaining obstruction would be entirely the
total-correlation term
\[
  \TC(X\vee Y)-\TC(X),
\]
which strengthens the case for connecting entropy loss with OR-collision and
Hall-defect structure.

\section{References}

\begin{thebibliography}{9}

\bibitem{Poonen}
B. Poonen,
\emph{Union-closed families},
Journal of Combinatorial Theory, Series A 59(2), 253--268, 1992.

\bibitem{BruhnSchaudt}
H. Bruhn and O. Schaudt,
\emph{The journey of the union-closed sets conjecture},
Graphs and Combinatorics 31, 2043--2074, 2015.

\bibitem{Gilmer}
J. Gilmer,
\emph{A constant lower bound for the union-closed sets conjecture},
arXiv:2211.09055, 2022.

\bibitem{AlweissHuangSellke}
R. Alweiss, B. Huang, and M. Sellke,
\emph{Improved lower bound for the union-closed sets conjecture},
Electronic Journal of Combinatorics 31(3), P3.35, 2024.

\bibitem{VuckovicZivkovic}
B. Vuckovic and M. Zivkovic,
\emph{Formalizing Frankl's Conjecture: FC-families},
arXiv:1207.3604, 2012.

\end{thebibliography}

\end{document}
